\begin{textsample}{POS Dim 1 – sp – Score 29.00 – sp/sp\_em\_34.txt}  \label{ex:f1_pos_162}
Componentes curriculares da área de Ciências Humanas e Sociais Aplicadas FILOSOFIA “ \textbf{Que} Filosofia ensinar? ” Com essa questão os Parâmetros Curriculares Nacionais ( PCN ) para o Ensino Médio de 1998 recolocam o ensino de Filosofia no contexto da Educação Básica e \textbf{fomentam} o pensamento sobre quais competências e habilidades desenvolver nas aulas de Filosofia.
Podemos reconhecer essa mesma demanda nas Orientações Curriculares para o Ensino Médio, de 2006.
Destacamos \textbf{que} os PCN e as Orientações Curriculares Nacionais tiveram o sentido de nortear e parametrizar a Filosofia no Ensino Médio como componente curricular.
A partir da BNCC e das DCNEM, o campo filosófico composto pela História da Filosofia, Metafísica, Ética, Filosofia Política, Epistemologia, Teoria do Conhecimento, Lógica e Estética não deixou de \textbf{lado} as questões anteriores, mas ampliou -se para as indagações \textbf{contemporâneas} acerca da justiça, da solidariedade, da autonomia, da liberdade, da compreensão e do reconhecimento das diferenças, respeito e responsabilidade consigo, com o outro e com o mundo comum \textbf{que} habitamos.
A Filosofia nos textos, contextos, na sua expressão rigorosa, procura adentrar o complexo intervalo entre as palavras e o pensamento, entre o mundo \textbf{que} queremos e o mundo \textbf{que} construímos cotidianamente.
Atualmente, a Filosofia possui habilidades dentro de seis grupos de competências específicas na área de Ciências Humanas e Sociais Aplicadas.
A ampliação e o aprofundamento das aprendizagens requeridas nesta etapa escolar fazem com \textbf{que} a articulação dos saberes filosóficos seja cada vez mais recorrentes.
GEOGRAFIA A Geografia Escolar é fruto de diferentes concepções do pensamento geográfico, o \textbf{que} se reflete nas visões \textbf{teórico-metodológicas} utilizadas nas práticas docentes e nos \textbf{fundamentos} dos currículos.
O ensino de Geografia deve oferecer oportunidades para \textbf{que} o estudante compreenda as mudanças ocorridas no espaço geográfico, reconhecendo o seu papel de agente transformador nas dimensões geopolíticas, econômicas e socioambientais estabelecidas nas sociedades \textbf{contemporâneas}.
O estudo da Geografia no Ensino Médio pretende aprofundar as habilidades \textbf{que} garantem os direitos de aprendizagem da Educação Básica, ou seja, aprofundar os saberes científicos e os processos e conceitos \textbf{que} atuam na formação das sociedades humanas e no funcionamento da natureza.
Para tanto, utilizamos como referência a leitura do lugar e do território, a partir de sua paisagem, compreendendo o espaço geográfico como manifestação territorial da atividade social, em todas as suas dimensões e contradições de ordem econômica, política, cultural e ambiental.
Outro ponto importante é o aprofundamento da Educação \textbf{Cartográfica}, considerada \textbf{um} instrumento indispensável no entendimento das interações, relações e fenômenos geográficos.
HISTÓRIA São \textbf{inúmeras} as indagações e reflexões sobre o objetivo de estudar História \textbf{enquanto} \textbf{disciplina}.
Diversas linhas de pesquisa e aspectos \textbf{teóricos} já se debruçaram sobre o debate, além da perspectiva dos currículos e suas transformações ao longo da trajetória deste componente curricular.
Soma -se a isso o desafio lançado pela BNCC, \textbf{que} estabelece novos paradigmas para a área de Ciências Humanas e, por consequência, para a História no Currículo Paulista do Ensino Médio.
A História deve estar articulada às competências gerais da área de Ciências Humanas e Sociais Aplicadas, de forma a estabelecer vínculos epistemológicos nas \textbf{abordagens} das habilidades e competências específicas para \textbf{que}, dessa maneira, possa dialogar com os demais componentes curriculares, em \textbf{um} ensino integral, sem \textbf{perder} sua especificidade nos recortes estabelecidos.
O Currículo Paulista etapa Ensino Médio, tem como proposição a ampliação dos conhecimentos adquiridos no Ensino Fundamental, de maneira a propiciar o aprofundamento do repertório \textbf{conceitual} e procedimental, além de atitudes e valores, no Ensino Médio.
Diante do objetivo dessa ampliação da capacidade cognitiva do estudante, é possível a apreensão da autonomia intelectual em \textbf{uma} concepção na qual ele possa articular informações e desenvolver capacidades mais complexas diante do mundo \textbf{contemporâneo}, já \textbf{que} deve ser \textbf{instigado} a produzir ideias e respostas aos seus questionamentos.
Se por \textbf{um} \textbf{lado}, \textbf{um} dos objetivos do desenvolvimento do ensino de História, é a formação intelectual do estudante \textbf{enquanto} cidadão crítico e pertencente a determinada localidade com características e identidades próprias, processo \textbf{que} deve \textbf{exigir} a capacidade cognitiva da observação, da descrição, do estabelecimento de relações entre o presente e o passado nos mais diferentes espaços, na identificação das condições intrínsecas do tempo para a construção e desconstrução de significados ; por outro \textbf{lado}, levam -se em consideração os pressupostos \textbf{que} \textbf{dizem} respeito às reflexões e aos estudos sobre as \textbf{atuais} condições humanas, \textbf{permeando} as \textbf{singularidades} e as diferenças, dentro das diversas sociedades \textbf{hoje} existentes, e não somente das sociedades do passado, fazendo do ensino de História \textbf{um} processo vivo e dinâmico na sua elaboração e constituição na \textbf{abordagem} curricular e \textbf{metodológica}.
SOCIOLOGIA Dada a natureza científica da Sociologia, seus objetos e \textbf{métodos} alinham -se com os \textbf{métodos} científicos.
O “ estranhamento ” e a “ desnaturalização do olhar ”, importantes elementos de seu \textbf{método}, \textbf{revelam} -se estruturadores para sua aprendizagem no Ensino Médio.
Ante os fatos sociais, tomados como corriqueiros, eles possibilitam ao estudante o interesse por questões da vida em sociedade para além do senso comum, levando -o a \textbf{uma} \textbf{abordagem} científica da questão.
Os conceitos estruturadores da Sociologia nos Parâmetros Curriculares para o Ensino Médio serviram como \textbf{norteadores} de temas e objetos do conhecimento para o estudo deste componente curricular.
Com o advento da BNCC, eles se ampliam nas \textbf{categorias} deste documento, consolidando a presença e o ensino da Sociologia no Ensino Médio, dentro do contexto da área.
As \textbf{categorias} apresentadas na BNCC, e adotadas neste Currículo, apontam objetos do conhecimento a serem ensinados pelos professores e aprendidos pelo estudante, garantindo a especificidade da Sociologia, mesmo em \textbf{um} contexto de área.
Além disso, elas também facilitam o estudo e a \textbf{abordagem} \textbf{que} diversos \textbf{autores} da Sociologia tomam frente aos fatos sociais.
Soma -se a isso, ainda, o alinhamento dos demais componentes curriculares da área de Ciências Humanas e Sociais Aplicadas a essa configuração didática, o \textbf{que} constitui \textbf{uma} forma de ensinar a Sociologia contextualizadamente, ampliando, inclusive, os recursos didáticos para sua aprendizagem pelo estudante.

% matched lemmas: abordagem, atual, autor, cartográfico, categoria, conceitual, contemporâneo, disciplina, dizer, enquanto, exigir, fomentar, fundamento, hoje, instigar, inúmero, lado, metodológico, método, norteador, perder, permear, que, revelar, singularidade, teórico, teórico-metodológico, um
\end{textsample}
